\section{A Perron--Frobenius structure, and what breaks it}\label{sec:cone}
\subsection{The cone the jump form carries}
Jump Dirichlet forms carry an order structure that quadratic forms in general do
not. By the Beurling--Deny criteria the semigroup and resolvent of such a form
are positivity preserving: if $u\geq0$ pointwise then $(\mathcal A+\lambda)^{-1}u\geq0$
for $\lambda>0$, where $\mathcal A$ is the generator. In a basis of indicator
cells on $I_L$, in which a nonnegative function is a nonnegative coefficient
vector, this says that the inverse matrix has no negative entries.

That is what happens. In a $48$-cell basis, the inverse of the archimedean
energy plus a positive constant has \emph{no} negative entries, at
$\lambda\in\{1,5\}$ and every $L$ tested; and so does the inverse of the
archimedean energy plus the constant $c_L$ wherever that operator is still
positive definite, which fails only at small $L$, where the bare archimedean form
is known to be indefinite.

\subsection{The pole form destroys it}
Adjoining the rank-two pole form $P_L$ removes the structure entirely.
\begin{center}
\small
\begin{tabular}{lrrrrr}
\toprule
$L$ & $4/5$ & $1$ & $3/2$ & $2$ & $3$\\
\midrule
negative fraction of $(K+5)^{-1}$
 & $0.000$ & $0.000$ & $0.000$ & $0.000$ & $0.000$\\
negative fraction of $(K+5+P_L)^{-1}$
 & $0.404$ & $0.511$ & $0.640$ & $0.703$ & $0.734$\\
\bottomrule
\end{tabular}
\end{center}
and the effect worsens with $L$. The mechanism is the Sherman--Morrison formula:
adding $\sigma ww^{\!\top}$ subtracts
$\sigma(A^{-1}w)(A^{-1}w)^{\!\top}/(1+\sigma w^{\!\top}A^{-1}w)$ from the
inverse, and when $A^{-1}\geq0$ entrywise and $w\geq0$ --- as $\cosh(x/2)>0$ is
--- that is a nonnegative rank-one subtraction, which drives entries negative.
\emph{Even the positive half of the pole form breaks inverse positivity.}

This is the third independent appearance of the pole term as the obstruction of
this subject, after the impossibility of supplying its negative direction by an
adjoined positive channel and the parity split \eqref{eq:parity-split}.

\subsection{What a cone can and cannot supply}
It is worth being precise, because the conclusion is not that cone methods are
unavailable but that three of the four standard ones cannot bound the quantity
in question.

\paragraph{(a) The existence of an invariant cone is automatic.} By
Vandergraft's theorem a matrix leaves some proper cone invariant as soon as its
spectral radius is an eigenvalue, which for a self-adjoint operator is always.
So the question is never whether a cone exists but whether a \emph{usable} one,
supplied by the structure, does.

\paragraph{(b) Krein--Rutman adds nothing to the spectral theorem here.} Its
content in dynamics is that a transfer operator is not self-adjoint and has no a
priori spectral theory, so the existence of a leading eigenpair is a real
statement. The operators of this subject are self-adjoint by construction and
their leading eigenpairs exist for free.

\paragraph{(c) Birkhoff's contraction bounds a gap, never a radius.} A positive
map taking a cone into itself with finite projective diameter $\Delta$ contracts
the Hilbert projective metric with ratio $\tanh(\Delta/4)$, which controls
$|\lambda_2|/|\lambda_1|$ and the convergence to the leading eigenvector. It
says nothing about $\lambda_1$, and Weil positivity is a statement about
$\lambda_1$ alone.

\paragraph{(d) Collatz--Wielandt does bound a radius, and needs a cone.} For a
cone-positive $M$,
\begin{equation}\label{eq:cw-abstract}
 \rho(M)=\inf\{\lambda:\exists\,u\in\operatorname{int}C,\ Mu\preceq\lambda u\},
\end{equation}
so a single pointwise supersolution certifies $\rho(M)\leq1$: it converts an
operator inequality into a pointwise inequality between two explicit functions.
That is what a cone is worth here, and Section~\ref{sec:cone}.2 is exactly what removes
it for $Q_L$.

The next section applies \eqref{eq:cw-abstract} to the largest piece of the
target that is still inside a cone.
